\section{Algoritmos}
\subsection{Reducción de base: Gram-Schmidt y LLL}
\begin{wrapfigure}{r}{5.2cm}
\vspace{-12pt}
\centering
% Base mala frente a base buena, con la celosia del Ejemplo 1 del enunciado.
%   (a) Base publica (mala):   (5, 19), (3, 14)
%   (b) Base reducida (buena): (3, 1), (-1, 4)   (LLL la devuelve con el signo cambiado)
% Las dos generan la MISMA celosia: los puntos son iguales en los dos dibujos.
% Lo unico que cambia es el paralelogramo de la base (sombreado): largo y estrecho
% con la base mala, casi un cuadrado con la buena.
\begin{tikzpicture}[x=0.14cm,y=0.14cm,
                    vec/.style={-{Stealth[length=4pt]},line width=0.9pt},
                    etiqueta/.style={fill=white,inner sep=0.8pt,font=\small}]
  \newcommand{\puntos}{%
    \foreach \a in {-8,...,10}{\foreach \b in {-3,...,6}{
      \fill[gray!70!black] ({3*\a-\b},{\a+4*\b}) circle (0.32);}}}

  % (a) Base publica
  \begin{scope}
    \clip (-3,-2) rectangle (12,22);
    \fill[gray!25] (0,0) -- (5,19) -- (8,33) -- (3,14) -- cycle;
    \puntos
    \draw[vec] (0,0) -- (5,19);
    \draw[vec] (0,0) -- (3,14);
  \end{scope}
  \draw[line width=0.4pt] (-3,-2) rectangle (12,22);
  \node[etiqueta,anchor=west] at (5.6,19) {$b_1$};
  \node[etiqueta,anchor=east] at (2.4,14) {$b_2$};
  \node[below,font=\small] at (4.5,-2) {(a) pública};

  % (b) Base reducida
  \begin{scope}[shift={(19,0)}]
    \begin{scope}
      \clip (-3,-2) rectangle (12,22);
      \fill[gray!25] (0,0) -- (3,1) -- (2,5) -- (-1,4) -- cycle;
      \puntos
      \draw[vec] (0,0) -- (3,1);
      \draw[vec] (0,0) -- (-1,4);
    \end{scope}
    \draw[line width=0.4pt] (-3,-2) rectangle (12,22);
    \node[etiqueta,anchor=west] at (3.5,0.5) {$b_1'$};
    \node[etiqueta,anchor=south] at (-1,4.6) {$b_2'$};
    \node[below,font=\small] at (4.5,-2) {(b) reducida};
  \end{scope}
\end{tikzpicture}

\caption{La celosía del Ejemplo~1 del enunciado con dos bases: los puntos son los mismos.}
\label{fig:celosia}
\vspace{-6pt}
\end{wrapfigure}
Una misma celosía tiene muchas bases. Con una buena, de vectores cortos y casi perpendiculares,
SVP y CVP son fáciles; con una mala, de vectores largos y casi paralelos, no
(Figura~\ref{fig:celosia}). GGH publica una base mala y guarda la buena, así que para atacarlo hay
que convertir la base pública en una buena. Eso es lo que hace LLL~\cite{lll,galbraith}.

LLL usa Gram-Schmidt para medir la base: los vectores perpendiculares $b_i^*$ (que no son de la
celosía) y los coeficientes $\mu_{ij}$, que indican cuánto se parece cada vector a los anteriores.
Con eso repite dos operaciones: restar a un vector un múltiplo entero de otro para que se parezcan
menos ($|\mu_{ij}|\le\frac12$), e intercambiar dos vectores seguidos cuando la parte perpendicular
del segundo, $b_k^*$, es bastante más corta que la del primero (condición de Lovász, con
$\delta=0{,}99$). Las dos son cambios de base unimodulares, como la $U$ de GGH, así que la celosía
no cambia, solo su descripción. LLL termina en tiempo polinómico y garantiza que el primer vector
es como mucho $(\delta-\frac14)^{-(n-1)/2}\approx1{,}16^{\,n-1}$ veces el más corto. Para medir la
calidad de una base usamos el cociente de Hadamard~\cite{galbraith}, $H=(\det L/\prod_i\|b_i\|)^{1/n}$,
que va de 0 a 1 y vale 1 si la base es perpendicular: en la tarea 3.03 pasa de 0{,}24 a 0{,}96.

\subsection{CVP aproximado: Babai, o por qué LLL rompe GGH}
El redondeo de Babai~\cite{babai} es justo el descifrado de GGH del enunciado: si $c=xB+e$,
entonces $cB^{-1}=x+eB^{-1}$, y al redondear sale $x$ si cada componente de $eB^{-1}$ está entre
$-\frac12$ y $\frac12$. Con la base pública $B^{-1}$ tiene números grandes y el redondeo falla; con
una base reducida $eB^{-1}$ es pequeño y acierta. Por eso no hace falta la clave privada: basta con
reducir la base pública y aplicar Babai. Usamos la variante del ``plano más cercano'', que redondea
vector a vector empezando por el último y acierta en más casos. Aun así, ninguna variante asegura
el punto más cercano (la traza 2.03 muestra un fallo), así que la usamos como punto de partida.

\subsection{Enumeración exacta (SVP y CVP)}
Para asegurar el óptimo hacemos una búsqueda exhaustiva, pero ordenada. Si $v=\sum_i x_ib_i$,
Gram-Schmidt permite escribir la distancia al objetivo como una suma de términos que nunca son
negativos~\cite{galbraith}:
\[
 \|v-t\|^2=\sum_{i=0}^{n-1}\|b_i^*\|^2\,(x_i-z_i)^2 ,
\]
donde el ``centro'' $z_i$ solo depende de $t$ y de los coeficientes $x_{i+1},\dots,x_{n-1}$. Por eso
elegimos los coeficientes de atrás hacia delante, como en un árbol. Partimos del punto de Babai
(en SVP, el vector más corto de la base reducida), cuya distancia $R^2$ sirve de límite. En cada
nivel solo probamos los enteros cercanos al centro que aún pueden quedar por debajo de $R^2$,
empezando por el más próximo (orden de Schnorr--Euchner~\cite{se}). Como lo que falta por sumar
no puede restar, en cuanto la suma parcial llega a $R^2$ abandonamos esa rama. Cada mejora baja
$R^2$, y cuando no quedan ramas sabemos que no existe un punto mejor. En SVP se descarta el cero;
las cuentas son con fracciones para no perder soluciones por redondeo. El coste es exponencial en
el peor caso, pero tras LLL el límite inicial es bueno y se visitan muy pocos nodos. Descartamos Lagrange--Gauss (solo vale en dimensión 2) y BKZ (un LLL más potente, innecesario
aquí): así un único método sirve para todas las tareas.

\subsection{LWE: fuerza bruta sobre el error y descifrado de un bit}
Si supiéramos el error de $n$ ecuaciones, despejaríamos $s$ como en un sistema normal (con la
inversa módulo $q$). Como los errores son pequeños, probamos todas las posibilidades en $[-c,c]$
para esas $n$ ecuaciones, despejamos $s$ y comprobamos si deja un error pequeño en \emph{todas}.
Empezamos con $c=0$ y vamos subiendo. Para cada $c$ hay $(2c+1)^n$ casos, que crecen muy rápido
con $n$; el mayor que usamos es L2.03, con $13^5\approx371\,000$, y paramos si se pasaría de un
millón. Las ecuaciones sobrantes son las que dan confianza: un secreto falso tendría que encajar en
todas por casualidad, y en la tarea L1.02 la probabilidad es de alrededor de $10^{-12}$. En cambio,
el Ejemplo~4 del enunciado solo tiene una ecuación de sobra, y con $c=1$ ya aparecen dos secretos
falsos, $(60,4,10,40)$ y $(6,19,60,37)$, antes que el verdadero.

Para descifrar un bit seguimos el enunciado: calculamos $R=C\cdot s \bmod q$ y miramos $d=B-R$
entre $-q/2$ y $q/2$. Con bit 0, $d$ es la suma de errores y queda cerca de 0; con bit 1 se le ha
sumado $\lfloor q/2\rfloor$ y queda cerca de $\pm q/2$. Decidimos con el umbral $|d|<q/4$.

\subsection{LWE como CVP: las mismas herramientas que en GGH}
\label{sec:lwecvp}
Para las LWE más grandes cambiamos de enfoque. Los vectores $As \bmod q$, para todos los $s$,
forman una celosía de $\mathbb Z^m$ (sumar dos de ellos, o sumar $q$ a una coordenada, da otro del
mismo tipo). Como $b=As+e$ con $e$ pequeño, $b$ está muy cerca del punto $As$: recuperar $s$ es un
CVP con objetivo $b$, igual que al descifrar GGH. Construimos una base de esa celosía con $A$ y
$q$, la reducimos con LLL, buscamos con la enumeración el punto más cercano a $b$ y de él
despejamos $s$. El coste ya no depende de $(2c+1)^n$ sino del tamaño de la celosía ($m\le30$ aquí). Comprobamos el resultado de tres formas: el error $b-As$ sale pequeño
(como mucho $\pm15$); L2.03, resuelta con los dos métodos, da el mismo secreto; y en todas menos
L3.02 el error mide menos de la mitad del vector más corto de la celosía, así que ningún otro punto
está igual de cerca. En L3.02 el margen no llega ($\|e\|^2=1073$ frente a $\lambda_1^2/4\approx1018$):
el punto es el más cercano, pero no podemos descartar del todo otro secreto con un error parecido.
Con pocas ecuaciones no funciona: en el Ejemplo~4 sale el secreto falso $(6,19,60,37)$.

\subsection{Ring-LWE: un sistema lineal por los parámetros de las tareas}
\label{sec:ring}
En Ring-LWE la clave pública es $b=a\cdot s+e$, con polinomios módulo $x^n+1$, coeficientes
módulo $q$ y $s$, $e$ pequeños~\cite{tema11}. Multiplicar por $a$ equivale a multiplicar por una
matriz $n\times n$ formada con sus coeficientes: es un LWE con estructura, cuya dificultad se apoya
en celosías ideales. Con parámetros bien elegidos no se conoce forma práctica de romperlo, y por
eso variantes suyas están en los estándares~\cite{lpr,fips203}.

Lo que hace vulnerables estas tareas son sus parámetros. Al mirar los datos vimos que $a$ solo
tiene coeficientes entre 0 y 15, cuando en un Ring-LWE normal ocupan todo el rango hasta $q-1$. El Tema~11 solo dice que $s$ y $e$ son
pequeños; nuestra hipótesis fue que están entre $-2$ y 2. Entonces cada coeficiente de $a\cdot s+e$
es una suma de $n$ productos de como mucho $15\cdot2$, y no pasa de $30n+2$: 482, 962 y 1922 para
$n=16$, 32 y 64, siempre por debajo de $q/2$ (699, 1559 y 2951). El módulo $q$ no llega a actuar y
$b$ es \emph{exactamente} $a\cdot s+e$ con enteros normales (que sin módulo el problema sea
lineal es conocido; lo nuestro fue verlo en estos datos). Resolvemos el sistema sin el error y redondeamos; que el redondeo acierte no lo demostramos, lo
comprobamos: ningún
valor se aleja más de 0{,}20 de un entero, el error $b-a\cdot s$ sale en $\{-1,0,1\}$ y $s$ en
$[-2,2]$, como suponíamos. Cuesta $O(n^3)$, frente a $5^{16}\approx10^{11}$ casos de una fuerza
bruta ya con $n=16$. Con $s$ desciframos como en el Tema~11: cada coeficiente de
$d=c_2-s\cdot c_1$ queda cerca de 0 (bit 0) o de $\pm q/2$ (bit 1). Con $a$ normal el módulo actúa
muchas veces y este atajo no sirve.
